\documentclass[11pt]{article}
\usepackage[a4paper,margin=22mm]{geometry}
\usepackage{fontspec}
\setmainfont{DejaVu Serif}
\setsansfont{DejaVu Sans}
\usepackage{microtype}
\usepackage{xcolor}
\usepackage{hyperref}
\usepackage{enumitem}
\usepackage{parskip}
\definecolor{Accent}{HTML}{971D32}
\definecolor{Muted}{HTML}{666666}
\hypersetup{colorlinks=true,urlcolor=Accent,linkcolor=Accent}
\setlist{leftmargin=1.6em,itemsep=0.3em,topsep=0.35em}
\setcounter{tocdepth}{2}
\renewcommand{\contentsname}{Contents}
\newcommand{\sectionrule}{\par\vspace{-0.5em}\noindent\textcolor{Accent}{\rule{\linewidth}{0.6pt}}\par}
\begin{document}

\begin{center}
{\sffamily\bfseries\color{Accent} TOEFL WORD OF THE DAY\par}
\vspace{8mm}
{\fontsize{42}{48}\selectfont\bfseries construe\par}
\vspace{2mm}
{\Large\sffamily\color{Accent} /kənˈstruː/\par}
\vspace{2mm}
{\small\sffamily Local audio phrase: construe\par}
{\small\sffamily Audio file: audios/construe.mp3\par}
\vspace{2mm}
{\large\sffamily\color{Muted} transitive verb\par}
\vspace{5mm}
{\sffamily interpret in a particular way\par}
\vspace{6mm}
{\large To construe something is to interpret what it means, especially when a statement, action, or piece of evidence could be understood in different ways.\par}
\vspace{6mm}
\begin{minipage}{0.84\textwidth}
\itshape “His silence could be construed as agreement with the proposal.”
\end{minipage}\par
\vspace{10mm}
\textbf{Date:} August 22nd 2026\quad\textbullet\quad \textbf{Level:} B1--mid-B2 TOEFL
\vfill
Created and compiled by \textbf{Amir Kooshky}\par
\vspace{2mm}
\href{https://t.me/KooshkyTOEFL}{Telegram Channel}\quad\textbullet\quad
\href{https://www.instagram.com/kooshkytoefl}{Instagram}\quad\textbullet\quad
\href{https://t.me/KooshkyTOEFL\_pv}{Personal Telegram}
\end{center}
\newpage
\tableofcontents
\newpage

\section{Meanings and usage}
\sectionrule
\subsection{Meaning 1: Interpret in a particular way}
\textbf{Part of speech:} transitive verb

\textbf{IPA:} /kənˈstruː/

\textbf{Pronunciation phrase:} construe

\textbf{Local audio file:} audios/construe.mp3

\textbf{Register:} formal; common in academic, legal, political, professional, and analytical writing

\textbf{Definition:} to understand or interpret the meaning of a statement, action, situation, or piece of information in a particular way

\textbf{In simpler words:} interpret or understand in a certain way

\textbf{Typical pattern:} construe + something + as + something

\subsubsection{Natural examples}
\begin{enumerate}
  \item His silence could be construed as agreement with the proposal.
  \item The researchers cautioned that the results should not be construed as proof of a causal relationship.
  \item Some employees construed the new policy as an attempt to reduce their independence.
  \item Her criticism was construed by some listeners as a personal attack.
  \item The law has traditionally been construed narrowly by the courts.
  \item A sudden decline in sales should not necessarily be construed as evidence that demand has permanently weakened.
  \item The committee construed the wording of the regulation differently from the previous administration.
  \item His decision to leave the meeting early was mistakenly construed as a sign of disrespect.
\end{enumerate}

\subsubsection{Nuance and implication}
\textbf{Central nuance:} Construe emphasizes the interpretation someone gives to words, behavior, evidence, or a situation rather than simply noticing what happened.

\textbf{Usual implication:} The thing being interpreted may allow more than one possible meaning, so different people may construe it differently.

\textbf{Compared with observe:} To observe something is to notice or examine it. To construe it is to decide what it means or what significance it has.

\subsubsection{Grammar notes}
\textbf{How to use it:} This verb normally needs the statement, action, evidence, or situation being interpreted directly after it.

\textbf{What can follow it:} The most important pattern is construe something as something, where as introduces the interpretation.

\textbf{Common preposition:} Use as after construe when stating the meaning given to something, as in construe silence as agreement.

\textbf{Position in a sentence:} The passive form is extremely common: something can be construed as something.

\subsubsection{Useful patterns}
\textbf{Pattern:} construe + something + as + something

\textbf{Use:} Use this when saying how a statement, action, or situation is interpreted.

\textbf{Example:} Some voters construed the announcement as a change in government policy.

\textbf{Pattern:} be construed as + noun or -ing form

\textbf{Use:} Use the passive form when the interpretation itself is more important than the person making it.

\textbf{Example:} The comment could be construed as criticism of the existing system.

\textbf{Pattern:} construe something broadly or narrowly

\textbf{Use:} Use this especially in legal or formal contexts when rules or language are given a wide or limited interpretation.

\textbf{Example:} The court construed the law narrowly.

\textbf{Pattern:} construe something differently

\textbf{Use:} Use this when people give different interpretations to the same information or behavior.

\textbf{Example:} Different readers may construe the final paragraph differently.

\textbf{Pattern:} should not be construed as

\textbf{Use:} Use this formal pattern to warn readers against a particular interpretation.

\textbf{Example:} These findings should not be construed as evidence that the treatment is effective in every patient.

\subsubsection{Common wrong usage}
\paragraph{Correction 1}
\textbf{Wrong:} The researchers construed the decline like evidence of failure.

\textbf{Correct:} The researchers construed the decline as evidence of failure.

\textbf{Why:} Use construe something as something, not like.

\paragraph{Correction 2}
\textbf{Wrong:} His silence was construed to agreement.

\textbf{Correct:} His silence was construed as agreement.

\textbf{Why:} Use as before the interpretation.

\paragraph{Correction 3}
\textbf{Wrong:} The statement can construe as a threat.

\textbf{Correct:} The statement can be construed as a threat.

\textbf{Why:} When the statement receives the interpretation, use the passive form be construed.

\paragraph{Correction 4}
\textbf{Wrong:} Construe means prove what something really means.

\textbf{Correct:} Construe means interpret what something means.

\textbf{Why:} A construal is an interpretation and may be disputed or incorrect; construe does not imply that the interpretation has been proved.

\section{Word family}
\sectionrule
\subsection{construal}
\textbf{Part of speech:} countable or uncountable noun

\textbf{IPA:} /kənˈstruːəl/

\textbf{Definition:} a particular way of understanding or interpreting something, or the act of interpreting it

\textbf{Relationship to the headword:} This noun names an interpretation or the process of construing something.

\textbf{Grammar pattern:} a construal of + statement, law, evidence, or situation

\textbf{Usage note:} It is formal and is especially common in legal, linguistic, philosophical, and academic writing.

\subsubsection{Examples}
\begin{enumerate}
  \item That construal of the law has been challenged by several legal scholars.
  \item Different construals of the same evidence can lead to very different conclusions.
  \item The court rejected the government's broad construal of the regulation.
\end{enumerate}

\subsection{misconstrue}
\textbf{Part of speech:} transitive verb

\textbf{IPA:} /ˌmɪskənˈstruː/

\textbf{Definition:} to interpret the meaning or intention of something incorrectly

\textbf{Relationship to the headword:} Misconstrue is formed by adding mis-, meaning wrongly, to construe.

\textbf{Grammar pattern:} misconstrue + something + as + something

\subsubsection{Examples}
\begin{enumerate}
  \item Please do not misconstrue my criticism as opposition to the entire project.
  \item The public may misconstrue the findings if the limitations are not explained clearly.
  \item His hesitation was misconstrued as a lack of interest.
\end{enumerate}

\section{Common collocations}
\sectionrule
\subsection{construe as}
\textbf{Applies to:} Interpret in a particular way

\textbf{Meaning:} interpret something as having a particular meaning

\textbf{Pattern:} construe + something + as + something

\textbf{Example:} The decision could be construed as a rejection of the earlier policy.

\subsection{be construed as}
\textbf{Applies to:} Interpret in a particular way

\textbf{Meaning:} be interpreted as having a particular meaning

\textbf{Pattern:} be construed as + noun or -ing form

\textbf{Example:} The statement was widely construed as criticism of the government.

\subsection{should not be construed as}
\textbf{Applies to:} Interpret in a particular way

\textbf{Meaning:} should not be interpreted to mean something

\textbf{Pattern:} should not be construed as + noun or -ing form

\textbf{Example:} The absence of evidence should not be construed as proof that no effect exists.

\subsection{could be construed as}
\textbf{Applies to:} Interpret in a particular way

\textbf{Meaning:} could reasonably be interpreted in a particular way

\textbf{Example:} His refusal to answer could be construed as an admission of uncertainty.

\subsection{construe broadly}
\textbf{Applies to:} Interpret in a particular way

\textbf{Meaning:} give words or rules a relatively wide interpretation

\textbf{Example:} The agency construed its authority broadly and introduced several new regulations.

\subsection{construe narrowly}
\textbf{Applies to:} Interpret in a particular way

\textbf{Meaning:} give words or rules a relatively limited interpretation

\textbf{Example:} The court construed the exception narrowly.

\subsection{construe differently}
\textbf{Applies to:} Interpret in a particular way

\textbf{Meaning:} give different meanings or interpretations to the same thing

\textbf{Example:} The two groups construed the same historical event differently.

\subsection{mistakenly construe}
\textbf{Applies to:} Interpret in a particular way

\textbf{Meaning:} interpret something incorrectly

\textbf{Example:} Readers may mistakenly construe uncertainty as evidence that the research has no value.

\subsection{a construal of}
\textbf{Applies to:} Interpret in a particular way

\textbf{Meaning:} a particular interpretation of something

\textbf{Pattern:} a construal of + law, statement, evidence, or event

\textbf{Example:} The judge rejected that construal of the statute.

\textbf{Usage note:} Construal is the noun form.

\section{Synonyms and antonyms}
\sectionrule
\subsection{Close synonyms}
\subsubsection{interpret}
\textbf{Part of speech:} transitive verb

\textbf{Applies to:} Interpret in a particular way

\textbf{Shared meaning:} Both mean understand something as having a particular meaning.

\textbf{Difference:} Interpret is the general and much more common word. Construe is more formal and is especially useful when discussing how words, actions, laws, or evidence are understood.

\textbf{Example:} Different historians interpret the event in different ways.

\subsubsection{understand}
\textbf{Part of speech:} transitive verb

\textbf{Applies to:} Interpret in a particular way

\textbf{Shared meaning:} Both can describe the meaning someone gives to words or behavior.

\textbf{Difference:} Understand is much broader and can simply mean know what something means. Construe emphasizes choosing or forming a particular interpretation.

\textbf{Example:} She understood his comment as a warning.

\subsubsection{read}
\textbf{Part of speech:} transitive verb

\textbf{Applies to:} Interpret in a particular way

\textbf{Shared meaning:} Both can mean interpret something as having a particular significance.

\textbf{Difference:} Read can informally mean interpret the significance of a situation or action. Construe is more formal and especially common in academic and legal contexts.

\textbf{Example:} Analysts read the announcement as a sign that the company was changing direction.

\subsection{Direct or useful antonyms}
\subsubsection{misconstrue}
\textbf{Part of speech:} transitive verb

\textbf{Applies to:} Interpret in a particular way

\textbf{Opposition:} Misconstrue means interpret something incorrectly, while construe itself is neutral about whether the interpretation is correct.

\textbf{Usage note:} This is not a perfect opposite because construe does not necessarily mean interpret correctly, but the contrast is educationally useful.

\textbf{Example:} Some readers misconstrued the author's criticism as a complete rejection of the theory.

\section{Words learners may confuse}
\sectionrule
\subsection{construct}
\textbf{Part of speech:} transitive verb

\textbf{Why learners confuse them:} The words look similar and share a historical root, but their modern meanings are different.

\textbf{Key difference:} Construe means interpret the meaning of something. Construct means build something physically or create something by putting parts together.

\textbf{construe:} The remark was construed as a criticism.

\textbf{construct:} Engineers constructed a bridge across the river.

\subsection{misconstrue}
\textbf{Part of speech:} transitive verb

\textbf{Why learners confuse them:} Misconstrue is formed directly from construe by adding the prefix mis-, meaning wrongly.

\textbf{Key difference:} Construe means interpret something in a particular way and is neutral about accuracy. Misconstrue specifically means interpret it incorrectly.

\textbf{construe:} The court construed the law as applying to digital records.

\textbf{misconstrue:} Some readers misconstrued the warning as a complete prohibition.

\subsection{infer}
\textbf{Part of speech:} transitive verb

\textbf{Why learners confuse them:} Both involve drawing meaning from information, but infer emphasizes reaching a conclusion while construe emphasizes interpretation.

\textbf{Key difference:} Construe means give a particular meaning to words, actions, or evidence. Infer means reach a conclusion from evidence or information that does not state the conclusion directly.

\textbf{construe:} She construed his silence as disagreement.

\textbf{infer:} From his silence, she inferred that he disagreed.

\section{Etymology}
\sectionrule
\textbf{Origin:} Construe ultimately comes from Latin construere, meaning to build or put together.

\textbf{Development:} The word developed from the idea of arranging or analyzing the parts of a sentence to the broader modern meaning of interpreting words, actions, or situations.

\textbf{Memory link:} When you construe something, you mentally put its pieces together to build an interpretation.

\section{Practice exercises}
\sectionrule
\subsection{Word family choice}
Choose the best member of the word family for each blank.

\begin{enumerate}
  \item The court rejected the government's broad \_\_\_\_\_ of the law.
  \item Readers may \_\_\_\_\_ the author's criticism as opposition to the entire proposal.
  \item The committee chose to \_\_\_\_\_ the regulation narrowly.
\end{enumerate}

\subsection{Correct or incorrect?}
Decide whether each sentence is correct. Correct the inaccurate ones.

\begin{enumerate}
  \item His silence was construed as agreement.
  \item The evidence should not be construed like proof of causation.
  \item Different courts may construe the same law differently.
  \item The statement can construe as a threat.
  \item To construe something means to prove its true meaning beyond doubt.
\end{enumerate}

\subsection{Collocation practice}
Complete each sentence with a natural collocation.

\begin{enumerate}
  \item The remark could be construed \_\_\_\_\_ criticism of the new policy.
  \item The findings should not be \_\_\_\_\_ as proof that the treatment works in every case.
  \item The court chose to construe the exception \_\_\_\_\_.
  \item Different observers may \_\_\_\_\_ the same behavior differently.
\end{enumerate}

\subsection{Complete answer key}
\subsubsection{Word family choice}
\begin{enumerate}
  \item \textbf{construal.} The court rejected the government's broad construal of the law. \emph{Use the noun construal for a particular interpretation.}
  \item \textbf{misconstrue.} Readers may misconstrue the author's criticism as opposition to the entire proposal. \emph{Use misconstrue when the interpretation is incorrect.}
  \item \textbf{construe.} The committee chose to construe the regulation narrowly. \emph{Use the base verb construe after to.}
\end{enumerate}

\subsubsection{Correct or incorrect?}
\begin{enumerate}
  \item \textbf{Correct} \emph{Be construed as is the standard passive pattern.}
  \item \textbf{Incorrect}. The evidence should not be construed as proof of causation. \emph{Use construe something as something.}
  \item \textbf{Correct} \emph{Construe something differently naturally describes different interpretations of the same text.}
  \item \textbf{Incorrect}. The statement can be construed as a threat. \emph{Use the passive form because the statement receives the interpretation.}
  \item \textbf{Incorrect}. To construe something means to interpret its meaning in a particular way. \emph{A construal is an interpretation and may be disputed or mistaken.}
\end{enumerate}

\subsubsection{Collocation practice}
\begin{enumerate}
  \item \textbf{as.} The remark could be construed as criticism of the new policy. \emph{Use construe something as something.}
  \item \textbf{construed.} The findings should not be construed as proof that the treatment works in every case. \emph{Should not be construed as is a common formal pattern for warning against an interpretation.}
  \item \textbf{narrowly.} The court chose to construe the exception narrowly. \emph{Construe narrowly means give a rule or statement a limited interpretation.}
  \item \textbf{construe.} Different observers may construe the same behavior differently. \emph{Construe differently means interpret the same thing in different ways.}
\end{enumerate}

\vfill
\begin{center}
\small\color{Muted} Created and compiled by \textbf{Amir Kooshky}\\
\href{https://t.me/KooshkyTOEFL}{Telegram Channel}\quad\textbullet\quad
\href{https://www.instagram.com/kooshkytoefl}{Instagram}\quad\textbullet\quad
\href{https://t.me/KooshkyTOEFL\_pv}{Personal Telegram}
\end{center}
\end{document}
